\documentclass[11pt]{article}
\usepackage[T1]{fontenc}
\usepackage[utf8]{inputenc}
\IfFileExists{lithuanian.ldf}{\usepackage[lithuanian]{babel}}{\usepackage{babel}\babelprovide[main,import]{lithuanian}}
\usepackage{lmodern,amsmath,amssymb,booktabs,array,tabularx,enumitem}
\usepackage[a4paper,margin=1.6cm]{geometry}
\setlength{\parindent}{0pt}
\setlength{\parskip}{5pt}
\setlength{\emergencystretch}{2em}
\setlist[itemize]{itemsep=3pt,topsep=4pt,parsep=0pt}
\newcommand{\sk}[1]{\par\medskip{\large\bfseries #1}\par}
\begin{document}
\begin{center}
{\Large\bfseries Kontrolinio darbo perrašymas}\\[4pt]
{\large Vilties variantas: sprendimai ir vertinimo instrukcija}
\end{center}
\textbf{Kvadratinės lygtys. 6 uždaviniai; 40 taškų; 45 min.}

\textbf{Galutinis taškų paskirstymas:} $6+6+8+6+8+6=40$.
Mokinio užduočių lape prie 2, 4 ir 5 uždavinių turi būti atitinkamai 6, 6 ir 8 taškai.

\sk{1 uždavinys -- 6 taškai}
\textbf{Sąlyga.} Išspręskite lygtį $(x+3)(x-2)=2x+6$.

\textbf{Sprendimas.}
\[
x^2+x-6=2x+6\quad\Longrightarrow\quad x^2-x-12=0.
\]
\[
D=(-1)^2-4\cdot1\cdot(-12)=49,\qquad
x=\frac{1\pm7}{2}.
\]
\textbf{Atsakymas:} $x=-3$ arba $x=4$.
\par\smallskip\renewcommand{\arraystretch}{1.2}\begin{tabularx}{\linewidth}{@{}X>{\centering\arraybackslash}p{1.1cm}@{}}\toprule Vertinamas pasiekimas & Taškai\\\midrule
Teisingai atskleidžia skliaustus: $x^2+x-6=2x+6$. & 2\\
Gauna kvadratinę lygtį $x^2-x-12=0$. & 1\\
Gauna $D=49$, išskaido $(x-4)(x+3)=0$ arba lygiaverčiai parengia lygtį šaknims rasti. & 1\\
Randa $x=4$. & 1\\
Randa $x=-3$. & 1\\
\textbf{Iš viso} & \textbf{6}\\
\bottomrule\end{tabularx}\par
Už skliaustų atskleidimą skirti 1 iš 2 taškų, jeigu teisingai sudarytos keturios sandaugos, bet padaryta viena ženklų ar skaičiavimo klaida.

\textbf{Alternatyvus sprendimas.}
\[
(x+3)(x-2)-2(x+3)=0\quad\Longrightarrow\quad(x+3)(x-4)=0.
\]
Toks sprendimas su abiem šaknimis vertinamas 6 taškais: 3 taškai už teisingą bendrojo daugiklio iškėlimą ir sandaugos lygtį, 1 taškas už nulinės sandaugos taisyklės pritaikymą ir po 1 tašką už kiekvieną šaknį. Taisyklė gali būti pritaikyta ir be atskiro žodinio paaiškinimo.

\textbf{Pastaba rengiant užduočių lapą.} Jei būtina tikrinti pertvarkymą į $ax^2+bx+c=0$, šį reikalavimą reikia aiškiai įrašyti į sąlygą. Dabartinė sąlyga leidžia trumpąjį sprendimą.

\newpage
\sk{2 uždavinys -- 6 taškai}
\textbf{Sąlyga.} Lygties $2x^2-6x+1=0$ šaknys yra $x_1$ ir $x_2$. Nerasdami šaknų, apskaičiuokite $x_1^3+x_2^3$.

\textbf{Sprendimas.} Pagal Vijetos teoremą
\[
x_1+x_2=3,\qquad x_1x_2=\frac12.
\]
\[
\begin{aligned}
x_1^3+x_2^3
&=(x_1+x_2)^3-3x_1x_2(x_1+x_2)\\
&=27-3\cdot\frac12\cdot3=\frac{45}{2}.
\end{aligned}
\]
\textbf{Atsakymas:} $\dfrac{45}{2}$.
\par\smallskip\renewcommand{\arraystretch}{1.2}\begin{tabularx}{\linewidth}{@{}X>{\centering\arraybackslash}p{1.1cm}@{}}\toprule Vertinamas pasiekimas & Taškai\\\midrule
Gauna $x_1+x_2=3$. & 1\\
Gauna $x_1x_2=\frac12$. & 1\\
Užrašo arba išveda tapatybę $x_1^3+x_2^3=(x_1+x_2)^3-3x_1x_2(x_1+x_2)$. & 2\\
Teisingai įstato reikšmes: $27-3\cdot\frac12\cdot3$. & 1\\
Apskaičiuoja $\frac{45}{2}$. & 1\\
\textbf{Iš viso} & \textbf{6}\\
\bottomrule\end{tabularx}\par
Už tapatybės žingsnį skirti 1 iš 2 taškų, jeigu mokinys teisingai išskleidžia $(x_1+x_2)^3$, bet neužbaigia reikiamo pertvarkymo.

\textbf{Kitas pilnavertis kelias.} Kiekviena šaknis tenkina $x_i^2=3x_i-\frac12$. Todėl
\[
x_i^3=3x_i^2-\frac12x_i=\frac{17}{2}x_i-\frac32,
\qquad
x_1^3+x_2^3=\frac{17}{2}\cdot3-3=\frac{45}{2}.
\]
Šiuo keliu skirti 1 tašką už šaknų sumą, 1 už ryšį $x_i^2=3x_i-\frac12$, 2 už kubo išraišką, 1 už abiejų kubų sumos sudarymą ir 1 už rezultatą.

\textbf{Jei mokinys pirmiausia randa šaknis,} už pilną teisingą tokį sprendimą skirti daugiausia 3 taškus: 1 už abi šaknis, 1 už teisingai sudarytą jų kubų sumą ir 1 už rezultatą. Neįvykdytas reikalavimas „nerasdami šaknų“.

\newpage
\sk{3 uždavinys -- 8 taškai}
\textbf{Sąlyga.} Raskite visas parametro $m$ reikšmes, su kuriomis lygties
\[
x^2-(m+2)x+2m=0
\]
realiosios šaknys tenkina $x_1^2+x_2^2=20$. Patikrinkite realiųjų šaknų egzistavimą.

\textbf{Sprendimas.} Pagal Vijetos teoremą
\[
x_1+x_2=m+2,\qquad x_1x_2=2m.
\]
\[
x_1^2+x_2^2=(x_1+x_2)^2-2x_1x_2=(m+2)^2-4m=m^2+4.
\]
\[
m^2+4=20\quad\Longrightarrow\quad m^2=16\quad\Longrightarrow\quad m=\pm4.
\]
Realiųjų šaknų egzistavimas:
\[
D=(m+2)^2-8m=(m-2)^2\ge0.
\]
Taigi abi parametro reikšmės tinka. \textbf{Atsakymas:} $m=-4$ arba $m=4$.
\par\smallskip\renewcommand{\arraystretch}{1.2}\begin{tabularx}{\linewidth}{@{}X>{\centering\arraybackslash}p{1.1cm}@{}}\toprule Vertinamas pasiekimas & Taškai\\\midrule
Gauna $x_1+x_2=m+2$. & 1\\
Gauna $x_1x_2=2m$. & 1\\
Pritaiko $x_1^2+x_2^2=(x_1+x_2)^2-2x_1x_2$. & 1\\
Sudaro $(m+2)^2-4m=20$. & 1\\
Teisingai suprastina iki $m^2=16$. & 1\\
Randa abi reikšmes $m=-4$ ir $m=4$. & 1\\
Pagrindžia realiųjų šaknų egzistavimą, kai $m=-4$. & 1\\
Pagrindžia realiųjų šaknų egzistavimą, kai $m=4$. & 1\\
\textbf{Iš viso} & \textbf{8}\\
\bottomrule\end{tabularx}\par
Už bendrą patikrą $D=(m-2)^2\ge0$ skiriami abu paskutiniai taškai. Vien už nepatikrintą teiginį „abi tinka“ šių taškų neskirti.

\textbf{Alternatyvus sprendimas.} Išskaidoma $(x-m)(x-2)=0$, todėl šaknys yra $m$ ir $2$. Gaunama $m^2+4=20$, taigi $m=\pm4$. Abi šaknys realios abiem gautais atvejais.

Šiuo keliu skirti 2 taškus už išskaidymą, 1 už abiejų šaknų nustatymą, 1 už sąlygos lygtį, 1 už $m^2=16$, 1 už abi parametro reikšmes ir 2 už realumo pagrindimą. Teisingai parodytos realiosios šaknys abiem atvejais yra pakankamas realumo pagrindimas; diskriminanto nereikalauti.

\newpage
\sk{4 uždavinys -- 6 taškai}
\textbf{Sąlyga.} Įvesdami tinkamą keitinį, išspręskite $x^4-5x^2-6=0$.

\textbf{Sprendimas.} Įvedame $t=x^2$, $t\ge0$:
\[
t^2-5t-6=0\quad\Longrightarrow\quad(t-6)(t+1)=0.
\]
Gauname $t=6$ arba $t=-1$. Antroji reikšmė netinka, nes $x^2\ge0$. Iš $x^2=6$ gauname $x=\pm\sqrt6$.

\textbf{Atsakymas:} $x=-\sqrt6$ arba $x=\sqrt6$.
\par\smallskip\renewcommand{\arraystretch}{1.2}\begin{tabularx}{\linewidth}{@{}X>{\centering\arraybackslash}p{1.1cm}@{}}\toprule Vertinamas pasiekimas & Taškai\\\midrule
Įveda keitinį $t=x^2$. & 1\\
Gauna $t^2-5t-6=0$. & 1\\
Randa $t=6$. & 1\\
Randa $t=-1$. & 1\\
Pagrįstai atmeta $t=-1$. & 1\\
Iš $x^2=6$ gauna abi reikšmes $x=\pm\sqrt6$. & 1\\
\textbf{Iš viso} & \textbf{6}\\
\bottomrule\end{tabularx}\par
Atmetimo pagrindimu laikomas ir iš anksto užrašytas $t\ge0$, jei toliau aiškiai pasirenkamas tik $t=6$. Jei užrašyta tik viena iš dviejų šaknų, paskutinis taškas neskiriamas.

Jei mokinys sprendžia be keitinio, išskaidydamas $(x^2-6)(x^2+1)=0$, už pilną teisingą sprendimą skirti daugiausia 4 taškus: 2 už išskaidymą, 1 už $x^2+1>0$ pagrindimą ir 1 už abi realiąsias šaknis. Taip išlaikomas sąlygoje reikalaujamo metodo vertinimas.

\sk{5 uždavinys -- 8 taškai}
\textbf{Sąlyga.} Įvesdami tinkamą keitinį, išspręskite
\[
\frac{x-2}{x+1}-2\frac{x+1}{x-2}=1.
\]
\textbf{Sprendimas.} $x\ne-1$, $x\ne2$. Įvedame
\[
t=\frac{x-2}{x+1},\qquad t\ne0.
\]
\[
t-\frac2t=1\quad\Longrightarrow\quad t^2-t-2=0
\quad\Longrightarrow\quad t=2\ \text{arba}\ t=-1.
\]
\[
\begin{aligned}
t=2:&\quad x-2=2(x+1)\quad\Longrightarrow\quad x=-4;\\
t=-1:&\quad x-2=-(x+1)\quad\Longrightarrow\quad x=\frac12.
\end{aligned}
\]
Abi reikšmės leidžiamos. \textbf{Atsakymas:} $x=-4$ arba $x=\frac12$.

\newpage
\sk{5 uždavinio vertinimas (tęsinys)}
\par\smallskip\renewcommand{\arraystretch}{1.2}\begin{tabularx}{\linewidth}{@{}X>{\centering\arraybackslash}p{1.1cm}@{}}\toprule Vertinamas pasiekimas & Taškai\\\midrule
Nurodo abu apribojimus $x\ne-1$, $x\ne2$ arba patikrina abu pradinius vardiklius gautoms šaknims. & 1\\
Įveda $t=\frac{x-2}{x+1}$ ir gauna $t-\frac2t=1$. & 1\\
Gauna $t^2-t-2=0$. & 1\\
Randa abi reikšmes $t=2$, $t=-1$. & 1\\
Kai $t=2$, gauna $x-2=2(x+1)$ arba lygiavertę tiesinę lygtį. & 1\\
Išsprendžia: $x=-4$. & 1\\
Kai $t=-1$, gauna $x-2=-(x+1)$ arba lygiavertę tiesinę lygtį. & 1\\
Išsprendžia: $x=\frac12$. & 1\\
\textbf{Iš viso} & \textbf{8}\\
\bottomrule\end{tabularx}\par
Atskiro užrašo $t\ne0$ nereikalauti, jei nurodyta $x\ne2$.

Jei lygtis išspręsta be keitinio, dauginant iš bendrojo vardiklio, skirti daugiausia 6 taškus: 1 už apribojimus arba patikrą, 1 už teisingą lygtį be vardiklių, 1 už kvadratinę lygtį $2x^2+7x-4=0$, 1 už tinkamą jos sprendimo eigą ir po 1 už kiekvieną šaknį.

\sk{6 uždavinys -- 6 taškai}
\textbf{Sąlyga.} Išskaidykite dauginamaisiais: a) $2x^2-7x+3$; b) $x^3-3x^2+2x-6$.

\textbf{a) Sprendimas.}
\[
D=49-24=25,\qquad x_1=\frac12,\quad x_2=3,
\qquad 2x^2-7x+3=(2x-1)(x-3).
\]
\par\smallskip\renewcommand{\arraystretch}{1.2}\begin{tabularx}{\linewidth}{@{}X>{\centering\arraybackslash}p{1.1cm}@{}}\toprule Vertinamas pasiekimas & Taškai\\\midrule
Randa vieną šaknį arba teisingai nustato vieną tiesinį daugiklį. & 1\\
Randa kitą šaknį arba teisingai nustato kitą tiesinį daugiklį. & 1\\
Užrašo tikslią sandaugą $(2x-1)(x-3)$, išlaikydamas vyriausiąjį koeficientą. & 1\\
\textbf{Iš viso už a)} & \textbf{3}\\
\bottomrule\end{tabularx}\par
Atsakymas $(x-\frac12)(x-3)$ vertinamas 2 taškais: šaknys nustatytos teisingai, bet prarastas koeficientas 2.

\textbf{b) Sprendimas.}
\[
x^3-3x^2+2x-6=x^2(x-3)+2(x-3)=(x-3)(x^2+2).
\]
\par\smallskip\renewcommand{\arraystretch}{1.2}\begin{tabularx}{\linewidth}{@{}X>{\centering\arraybackslash}p{1.1cm}@{}}\toprule Vertinamas pasiekimas & Taškai\\\midrule
Tikslingai sugrupuoja: $(x^3-3x^2)+(2x-6)$. & 1\\
Iškelia bendruosius daugiklius: $x^2(x-3)+2(x-3)$. & 1\\
Užbaigia: $(x-3)(x^2+2)$. & 1\\
\textbf{Iš viso už b)} & \textbf{3}\\
\bottomrule\end{tabularx}\par
Abiejose dalyse teisinga galutinė sandauga pakankamai parodo rezultatą ir vertinama visais 3 taškais. Tarpinių eilučių nereikalauti.

\newpage
\sk{Bendrosios vertinimo taisyklės}
\begin{enumerate}[leftmargin=*,itemsep=8pt]
\item \textbf{Vienos klaidos pakartotinai nebausti.} Po vienos aritmetinės ar ženklų klaidos tolesnius teisingus veiksmus vertinti pagal mokinio gautus duomenis, jei klaida iš esmės nesupaprastino uždavinio. Nauja nepriklausoma klaida vertinama atskirai.
\item \textbf{Priimti lygiaverčius metodus.} Kai sąlyga aiškiai reikalauja konkretaus būdo, taikyti prie atitinkamo uždavinio nurodytas ribas. Trumpesnis pagrįstas sprendimas savaime nėra mažiau vertas.
\item \textbf{Tarpiniai žingsniai gali būti numanomi iš teisingo pertvarkymo.} Nereikalauti atskiros eilutės vien tam, kad būtų pažymėtas kriterijus. Pavyzdžiui, tinkama sandauga ir abi šaknys gali pagrįsti kelis susijusius pasiekimus.
\item \textbf{Už atsakymo perrašymą į langelį papildomai taškų neskirti ir neatimti.} Pakanka, kad rezultatas būtų aiškiai matomas sprendime. Vertinti galutinį mokinio pasirinktą atsakymą ir jo sprendimą.
\item \textbf{Neteisingas atsakymas nepanaikina teisingos eigos.} Skirti taškus už teisingai atliktus vertinamus žingsnius, atsižvelgiant į vienos klaidos taisyklę.
\item \textbf{Vien atsakymas nėra visas sprendimas.} Jei jokios eigos nėra, skirti tik lentelėje numatytus galutinio rezultato taškus: 1 uždavinyje iki 2; 2 uždavinyje 1; 3 uždavinyje 1; 4 uždavinyje 1; 5 uždavinyje iki 2. 6 uždavinyje teisinga sandauga vertinama visais atitinkamos dalies taškais.
\item \textbf{Vertinti sveikaisiais taškais.} Netaikyti papildomų bendrų nuoskaitų už tvarką ar pasirinktą žymėjimą, jei matematinė mintis aiški.
\item \textbf{Nedubliuoti vertinimo.} Alternatyvių sprendimo kelių taškai nesumuojami. Už vieną uždavinį negalima skirti daugiau nei jo maksimalus taškų skaičius.
\item \textbf{Pažymio apvalinimas.} Surinktų taškų procentinę išraišką padalyti iš 10. Kai surinkta mažiau nei 70\%, gautą skaičių apvalinti į didesnę pusę iki sveikojo pažymio: 61\% atitinka 6,1, todėl pažymys yra 7. Nuo 70\% taikyti įprastą matematinį apvalinimą: 72\% atitinka 7,2, todėl pažymys yra 7, o 76\% atitinka 7,6, todėl pažymys yra 8.
\end{enumerate}
\medskip
\begin{center}
\begin{tabular}{lrrrrrrr}
\toprule
Uždavinys & 1 & 2 & 3 & 4 & 5 & 6 & Iš viso\\
\midrule
Taškai & 6 & 6 & 8 & 6 & 8 & 6 & \textbf{40}\\
\bottomrule
\end{tabular}
\end{center}
\end{document}